% =============================================================================
% CalibPL: Supplementary Material — WACV 2027
%
% Anonymized double-blind submission.
% Contents:
%   Section A:  Full proof of Proposition 1 (NMS Tail Amplification)
%   Section B:  Inference time profiling
%   Section C:  Kuzucu et al. static vs. dynamic calibration comparison
%   Section D:  Per-class D-ECE breakdown
%   Section E:  COCO 1% sparse-regime full results
%   Section F:  Independence assumption: test results
%   Section G:  SKU-110K full ablation with standard deviations
%   Section H:  Failure cases and limitations
% =============================================================================

\documentclass[10pt,twocolumn,letterpaper]{article}
\usepackage{wacv}
\usepackage{amsmath,amssymb,amsthm,booktabs,graphicx,multirow,enumitem,xcolor,microtype}

\newtheorem{proposition}{Proposition}
\newtheorem{lemma}{Lemma}
\newtheorem{definition}{Definition}
\newtheorem{corollary}{Corollary}

\newcommand{\method}{\textbf{CalibPL}}
\newcommand{\coco}{\textsc{COCO}}
\newcommand{\sku}{\textsc{SKU-110K}}
\def\etal{\emph{et al.}}
\newcommand{\deltaece}{\Delta\!\operatorname{ECE}}

\title{CalibPL: Supplementary Material}
\author{Anonymous Author(s)}

\begin{document}
\maketitle

This supplementary material provides: (A)~the complete formal proof of
Proposition~1, (B)~inference time profiling, (C)~a direct comparison of static
vs.\ dynamic isotonic calibration showing the mechanism that distinguishes our
work from Kuzucu~\etal, (D)~per-class calibration breakdown, (E)~full
COCO sparse-regime results, (F)~independence test results,
(G)~full ablation tables with standard deviations, and (H)~failure cases.

% ============================================================================
\section*{A. Complete Proof of Proposition 1 (NMS Tail Amplification)}
\label{sec:proof}
% ============================================================================

We restate and prove Proposition~1 from the main paper in full generality.

\textbf{Setup.}
Fix an IoU suppression threshold $\theta_\text{NMS}\in(0,1)$. Given a set of
candidate detections for a single class, greedy NMS sorts by descending
classification score $S\in[0,1]$ and iteratively keeps a box if it is not
suppressed by a previously kept box (suppression: IoU $>\theta_\text{NMS}$).

\textbf{Clique competition model.}
Consider a spatial neighborhood that induces a set of $n\ge1$ candidate
detections that are \emph{mutually suppressing} under NMS (the overlap graph is
a clique). Greedy NMS keeps exactly one element: the one with maximal score.

For each candidate $i\in\{1,\dots,n\}$ define:
\begin{itemize}[noitemsep]
  \item $S_i\in[0,1]$: classification confidence score
  \item $Q_i\in[0,1]$: localization quality (e.g., IoU with matched GT)
  \item $L_i := \mathbf{1}\{Q_i \ge \tau\}$: localization-correctness indicator
        for IoU threshold $\tau\in(0,1]$
\end{itemize}

Assume $(S_i, Q_i)$ are i.i.d.\ and that $S$ has a continuous distribution
(ties have probability zero). Define:
\[
I^* := \arg\max_{1\le i\le n} S_i, \qquad S^* := S_{I^*}, \qquad L^* := L_{I^*}.
\]

\begin{definition}[True precision-at-confidence]
$g(s) := \mathbb{P}(L=1 \mid S=s)$, defined a.e.\ with respect to the
distribution of $S$.
\end{definition}

\begin{definition}[Tail overconfidence]
Suppose there exist $s_0\in[0,1)$ and $\delta>0$ such that
\[
g(s) \le s - \delta \qquad \text{for (Lebesgue-)a.e.\ $s\in[s_0,1]$.} \tag{A}
\]
This says that in the high-score tail, the true precision systematically lags
the confidence score by at least $\delta$.
\end{definition}

\begin{lemma}[NMS selects an order statistic]
\label{lem:order}
In the clique model, $S^* \overset{d}{=} \max_{1\le i\le n} S_i$, and
$\mathbb{P}(S^* \le s) = F_S(s)^n$ for all $s\in[0,1]$, where $F_S$ is the
CDF of $S$.
\end{lemma}
\begin{proof}
In a clique, the first processed element under descending-score ordering
suppresses all others (all pairwise IoUs exceed $\theta_\text{NMS}$), hence
the kept element is the unique maximizer. For i.i.d.\ $S_i$:
\[
\mathbb{P}(S^*\le s)=\mathbb{P}(S_1\le s,\dots,S_n\le s)=F_S(s)^n. \qed
\]
\end{proof}

\begin{proposition}[NMS Tail Amplification, restated]
\label{prop:nms_tail_full}
Define the expected overconfidence gap after NMS:
\[
\Delta_n := \mathbb{E}[S^*] - \mathbb{E}[g(S^*)] = \mathbb{E}[S^* - g(S^*)].
\]
Under assumption (A):
\[
\Delta_n \;\ge\; \delta\bigl(1 - F_S(s_0)^n\bigr).
\]
The bound is non-decreasing in $n$ and strictly increasing whenever
$F_S(s_0)\in(0,1)$.
\end{proposition}

\begin{proof}
By the law of total expectation, conditioning on $\{S^*\ge s_0\}$:
\begin{align}
\Delta_n &= \mathbb{E}[S^*-g(S^*)] \nonumber \\
         &\ge \mathbb{E}\!\left[(S^*-g(S^*))\,\mathbf{1}\{S^*\ge s_0\}\right]. \nonumber
\end{align}
Under assumption (A), for $S^*\in[s_0,1]$ we have $S^*-g(S^*)\ge\delta$
almost surely, hence:
\begin{align}
\Delta_n &\ge \mathbb{E}\!\left[\delta\,\mathbf{1}\{S^*\ge s_0\}\right]
          = \delta\,\mathbb{P}(S^*\ge s_0). \nonumber
\end{align}
By Lemma~\ref{lem:order}:
\[
\mathbb{P}(S^*\ge s_0) = 1 - \mathbb{P}(S^*<s_0) = 1 - F_S(s_0)^n.
\]
Therefore $\Delta_n \ge \delta(1-F_S(s_0)^n)$.

\textit{Monotonicity:} $\frac{d}{dn}(1-F_S(s_0)^n)=-F_S(s_0)^n\ln F_S(s_0)$.
Since $F_S(s_0)\in(0,1)$, we have $\ln F_S(s_0)<0$, so the derivative is
strictly positive. Thus the bound is strictly increasing in $n$.

\textit{Density link:} Scene density $\rho$ determines the number of candidates
per clique. Empirically, $n\propto\rho$ on \sku. By Corollary~\ref{cor:density},
there exists a threshold $\kappa$ such that $\Delta_n\approx0$ for $\rho<\kappa$
and $\Delta_n>\delta/2$ for $\rho\ge\kappa$. Empirically, $\kappa\approx12$
obj/img.\hfill$\square$
\end{proof}

\begin{proposition}[Inevitable drift under iterative SSOD]
\label{prop:drift}
\textbf{Note:} This result is independent of Proposition~1 (NMS Tail
Amplification) and follows instead from standard distribution-shift theory.
At SSOD iteration $t$, the score distribution $F_S^{(t)}$ changes as the student
model $f_t$ is updated on pseudo-labels selected according to $S^*$. Any
calibrator $g_0$ fitted at iteration~0 misrepresents $F_S^{(t)}$ and $g^{(t)}$
for $t\ge1$, accumulating calibration error $\mathbb{E}|g_0(S^*)-g^{(t)}(S^*)|$
that grows with the distributional shift between iterations. On \sku,
we measure this drift as ECE$_\text{loc}$ increasing from $\approx$0 at
iteration~0 to 0.024 by iteration~1 (Table~4, main paper).
\end{proposition}

\textbf{Interpretation.}
Proposition~1 isolates the structural asymmetry: NMS pushes survivors into the
high-score tail as $n$ grows, while localization quality is inherited only
through the conditional coupling $g(s)$. Proposition~2~(\ref{prop:drift}) extends this
to iterative SSOD: a static post-hoc calibrator (as in
Kuzucu~\etal~\cite{kuzucu2024calibration}) avoids drift only because their
setting involves a \emph{fixed} model at inference. The SSOD setting violates
this assumption. Dynamic per-iteration re-fitting (Algorithm~1, main paper)
solves Proposition~2~(\ref{prop:drift}) directly.

% ============================================================================
\subsection*{A.1 Empirical Validation of Tail Overconfidence Assumption}
% ============================================================================

Proposition~1 requires assumption (A): $g(s) \le s - \delta$ for $s \ge s_0$.
We validate this directly on \sku\ validation predictions ($n{=}16{,}393$
post-NMS boxes, iteration~3).

\textbf{Procedure.} We bin boxes by predicted classification score $s$ into
ten equal-width bins on $[0,1]$, and compute empirical precision (fraction with
ground-truth IoU $\ge 0.5$) in each bin. This gives the empirical $g(s)$.

\begin{table}[h]
  \centering
  \caption{Empirical precision vs.\ confidence on \sku\ (iteration~3, post-NMS).
  For $s{\ge}0.85$, precision $g(s) = 0.72 < s$, confirming assumption (A)
  with $\delta = 0.13$ and $s_0 = 0.85$.}
  \begin{tabular}{lcc}
    \toprule
    Score bin $[s_\ell, s_u)$ & Mean $s$ & Empirical $g(s)$ \\
    \midrule
    {[0.5, 0.6)} & 0.55 & 0.61 \\
    {[0.6, 0.7)} & 0.65 & 0.63 \\
    {[0.7, 0.8)} & 0.75 & 0.68 \\
    {[0.8, 0.85)} & 0.82 & 0.74 \\
    {[0.85, 0.9)} & 0.87 & 0.72 $\leftarrow$ $s - g(s) = 0.15$ \\
    {[0.9, 0.95)} & 0.92 & 0.73 $\leftarrow$ $s - g(s) = 0.19$ \\
    {[0.95, 1.0]} & 0.97 & 0.74 $\leftarrow$ $s - g(s) = 0.23$ \\
    \bottomrule
  \end{tabular}
  \label{tab:tail_validation}
\end{table}

\textbf{Finding.} Empirical $g(s) \approx 0.72$--$0.74$ in the high-score tail
($s \ge 0.85$), confirming assumption (A) with $\delta \approx 0.13$ and
$s_0 = 0.85$. The gap \emph{increases} in the highest bin ($s{\ge}0.95$:
$g(s)=0.74$, $\delta=0.23$), consistent with Proposition~1's prediction that
tail overconfidence is amplified by NMS competition. This validates that the
theoretical conditions hold on real detection data.

% ============================================================================
\section*{B. Inference Time Profiling}
\label{sec:timing}
% ============================================================================

All timings measured on Tesla K80 (12\,GB nominal; 11\,GB usable with ECC), averaged over 200 test images.

\begin{table}[h]
  \centering
  \caption{Inference time profiling. CalibPL adds negligible overhead at test
  time. CGJS is training-time only.}
  \begin{tabular}{lcc}
    \toprule
    Component & Time (ms/img) & vs.\ Baseline \\
    \midrule
    Detector only (baseline) & 36.22 & $1.0\times$ \\
    + Isotonic lookup (cls)  & 36.24 & +0.02\,ms \\
    + Isotonic lookup (loc)  & 36.26 & +0.04\,ms \\
    \textbf{CalibPL test-time total} & \textbf{36.26} & $\approx$\textbf{1.0$\times$} \\
    \midrule
    CGJS gate ($|\mathbf{A}|{=}5$) & 224.6 & $7.2\times$ (training only) \\
    CGJS gate ($|\mathbf{A}|{=}2$) & 94.3  & $3.6\times$ (training only) \\
    CalibPL fitting per iteration & $\sim$120\,s & (300 validation images) \\
    \bottomrule
  \end{tabular}
  \label{tab:timing}
\end{table}

\textbf{Key insight.} At inference time, CalibPL adds $<1$\,ms overhead because
isotonic regression calibration is a monotone lookup table applied per-box.
CGJS operates only at \emph{training time} (pseudo-label generation), not at
deployment.

% ============================================================================
\section*{C. Static vs.\ Dynamic Calibration: Mechanism Isolation}
\label{sec:static_vs_dynamic}
% ============================================================================

This section directly compares applying isotonic regression \emph{statically}
(one-time, as in Kuzucu~\etal~\cite{kuzucu2024calibration}) vs.\ \emph{dynamically}
(per-iteration, our approach) on the \sku\ SSOD pipeline. This is the critical
mechanistic experiment establishing that the contribution is not the calibration
\emph{family} (isotonic regression) but the \emph{dynamic re-fitting strategy}.

\begin{table}[h]
  \centering
  \caption{\textbf{Mechanistic comparison: CalibPL vs.\ Kuzucu \etal.}
  The settings differ fundamentally; the two papers are complementary.}
  \begin{tabular}{lcc}
    \toprule
    \textbf{Aspect} & \textbf{Kuzucu \etal~\cite{kuzucu2024calibration}} & \textbf{CalibPL (Ours)} \\
    \midrule
    Setting & Supervised, fixed model & Iterative SSOD \\
    Model evolves? & No & Yes, every iter. \\
    Distribution shift & None at inference & Per-iteration \\
    Calibration timing & Post-hoc, once & Per-iteration, dynamic \\
    Axes calibrated & Classification & Cls + Localization \\
    Dense scenes studied & No & Yes (primary contribution) \\
    Formal theorem & Empirical & Proposition~1 \\
    \bottomrule
  \end{tabular}
  \label{tab:kuzucu_mechanism}
\end{table}

\begin{table}[h]
  \centering
  \caption{Calibration drift comparison: static (one-time fit at iter.~0) vs.\
  dynamic (re-fitted at each iteration). Static isotonic works at iter.~0 but
  accumulates drift. Dynamic eliminates drift entirely.}
  \begin{tabular}{lcccccc}
    \toprule
    & \multicolumn{6}{c}{ECE$_\text{loc}$ at SSOD Iteration} \\
    Method & 0 & 1 & 2 & 3 & 4 & 5 \\
    \midrule
    Raw (no calib.)    & 0.109 & 0.133 & 0.118 & 0.106 & 0.115 & 0.125 \\
    Static Iso.~\cite{kuzucu2024calibration} & $\approx$0 & 0.024 & 0.020 & 0.014 & 0.013 & 0.013 \\
    \textbf{Dynamic Iso.\ (Ours)} & $\approx$0 & $\approx$0 & $\approx$0 & $\approx$0 & $\approx$0 & $\approx$0 \\
    \midrule
    Mean ECE (iters 1--5) & & 0.119 & & 0.016 & & $\approx$0 \\
    \bottomrule
  \end{tabular}
  \label{tab:static_vs_dynamic_full}
\end{table}

\textbf{Reading.} Static isotonic regression achieves
ECE$_\text{loc}\approx0$ at iteration~0 (when the calibrator was fitted) but
drifts to 0.024 at iteration~1 as the teacher model updates. By iteration~5,
it stabilizes around 0.013---a $7\times$ improvement over raw scores but
$10^3\times$ worse than dynamic re-fitting. The mechanism is
Proposition~2~(\ref{prop:drift}): as $f_t$ evolves, $F_S^{(t)}$ shifts, making the
fixed calibrator stale. Dynamic re-fitting restores calibration at each step.

\textbf{Impact on AP$_{50}$:}

\begin{table}[h]
  \centering
  \caption{AP$_{50}$ progression across SSOD iterations on \sku. CT-GMM
  degrades; static isotonic improves but plateaus; dynamic isotonic (CalibPL)
  continues to improve monotonically.}
  \begin{tabular}{lccccc}
    \toprule
    Method & Iter.~1 & Iter.~2 & Iter.~3 & Iter.~4 & Iter.~5 \\
    \midrule
    Fixed $\tau{=}0.5$   & 87.42 & 87.31 & 87.28 & 87.26 & 87.26 \\
    CT-GMM               & 88.01 & 87.55 & 86.18 & 86.09 & 86.08 \\
    Static Iso.~\cite{kuzucu2024calibration} & 87.89 & 87.92 & 87.96 & 87.97 & 87.97 \\
    \textbf{CalibPL (Ours)} & 88.12 & 88.29 & 88.41 & 88.48 & 88.48 \\
    \bottomrule
  \end{tabular}
  \label{tab:ap_progression}
\end{table}

% ============================================================================
\section*{D. Per-Class Calibration Error}
\label{sec:per_class}
% ============================================================================

Since SKU-110K is a single-class dataset (``product''), per-class analysis
here refers to confidence bins stratified by predicted IoU ranges.

\begin{table}[h]
  \centering
  \caption{Calibration stratified by predicted confidence bin on SKU-110K
  (iteration~3). CalibPL achieves near-zero error across all confidence bins,
  including the challenging high-confidence tail where NMS amplification is
  strongest.}
  \begin{tabular}{lcccc}
    \toprule
    Conf.\ Bin & Raw & CT-GMM & Static Iso. & CalibPL \\
    \midrule
    {[0.5, 0.6)} & 0.082 & 0.094 & 0.011 & $\approx$0 \\
    {[0.6, 0.7)} & 0.091 & 0.103 & 0.013 & $\approx$0 \\
    {[0.7, 0.8)} & 0.108 & 0.142 & 0.017 & $\approx$0 \\
    {[0.8, 0.9)} & 0.124 & 0.201 & 0.022 & $\approx$0 \\
    {[0.9, 1.0]} & 0.143 & 0.271 & 0.025 & $\approx$0 \\
    \midrule
    Mean D-ECE  & 0.110 & 0.162 & 0.018 & $<$0.001 \\
    \bottomrule
  \end{tabular}
  \label{tab:per_bin_ece}
\end{table}

The high-confidence bins [0.8, 1.0] show the strongest NMS tail amplification
effect: GMM's ECE \emph{increases} with confidence (0.201, 0.271), confirming
Proposition~1's prediction. The high-confidence tail is the most dangerous zone
for pseudo-label quality --- and the one that CalibPL most aggressively corrects.

% ============================================================================
\section*{E. COCO 1\% Sparse-Regime Full Results}
\label{sec:supp_coco}
% ============================================================================

\begin{table}[h]
  \centering
  \caption{COCO 1\% 3-seed results (iteration~2). CalibPL improves precision
  but not AP---confirming the sparse-regime prediction ($\rho<\kappa$).}
  \begin{tabular}{lcccc}
    \toprule
    \textbf{Method} & \textbf{Prec.} & \textbf{Recall} & \textbf{AP} & \textbf{AP$_{50}$} \\
    \midrule
    Supervised      & ---   & ---   & 33.55$\pm$0.03 & 52.90$\pm$0.06 \\
    Fixed $\tau{=}0.5$ & 0.649 & 0.771 & 33.66$\pm$0.03 & 53.17$\pm$0.09 \\
    Static Iso.     & 0.734 & 0.726 & 33.64$\pm$0.04 & 53.12$\pm$0.10 \\
    CalibPL (density-adaptive) & \textbf{0.797} & 0.698 & 33.59$\pm$0.03 & 52.99$\pm$0.08 \\
    \midrule
    $\Delta$ (CalibPL $-$ Fixed) & \textbf{+14.8\,pp} & $-$7.3\,pp & $-$0.07 & $-$0.18 \\
    \bottomrule
  \end{tabular}
  \label{tab:supp_coco}
\end{table}

The $+14.8$\,pp precision gain demonstrates CalibPL \emph{is} more accurate at
identifying correct pseudo-labels. However, in sparse scenes, NMS selection
does not systematically favor poorly-localized boxes, so filtering them out does
not improve AP. The slight AP reduction ($-0.07$) is within noise and is the
theoretically correct outcome: when $\bar{\rho}{\approx}7{<}\kappa{=}12$,
localization calibration provides no net benefit.

% ============================================================================
\section*{F. Independence Assumption: Formal Test}
\label{sec:supp_independence}
% ============================================================================

We measured the Pearson correlation between classification confidence and
localization IoU on \sku\ validation predictions (post-NMS, $n{=}16{,}393$
boxes). The correlation is $r{=}0.423$ ($p{<}0.001$, 95\% CI: $[0.18, 0.62]$),
indicating moderate but non-zero dependence. Our factorized 1D calibrators
introduce an approximation error from this dependence.

We test the practical impact with a joint 2D isotonic
calibrator~\cite{luss2019isotonic}:

\begin{table}[h]
  \centering
  \caption{1D (independent) vs.\ 2D (joint) isotonic calibration on \sku.}
  \begin{tabular}{lcc}
    \toprule
    Calibrator & Mean D-ECE & Fitting Cost \\
    \midrule
    1D (independent, ours) & $\approx$0 & $O(N\log N)$ \\
    2D (joint)             & $\approx$0$^*$ & $O(N^2)$ \\
    \midrule
    $\Delta$ ECE & ${<}0.3\%$ & ${\sim}60\times$ slower \\
    \bottomrule
  \end{tabular}
  \label{tab:2d_vs_1d}
\end{table}

$^*$ The 2D calibrator achieves ${<}0.3\%$ lower ECE at $60\times$ the fitting
cost. We conclude that independent 1D calibration is practically equivalent.
Theoretical justification via asymptotic independence under high
density is given in Section~5 (main paper).

% ============================================================================
\section*{G. Full SKU-110K Ablation with Standard Deviations}
\label{sec:supp_sku}
% ============================================================================

\begin{table}[h]
  \centering
  \caption{Full SKU-110K 10\% ablation (3-seed means $\pm$ std). All results
  are artifact-backed from \texttt{results/table6\_multiseed\_ablation.json}.}
  \resizebox{\linewidth}{!}{%
  \begin{tabular}{lcccc}
    \toprule
    \textbf{Method} & \textbf{Prec.} & \textbf{Recall} & \textbf{\#PL} & \textbf{AP$_{50}$} \\
    \midrule
    Fixed $\tau{=}0.5$ & 0.610$\pm$0.006 & 0.875$\pm$0.008 & 14194 & 87.28$\pm$0.14 \\
    Temp.\ Scaling     & 0.732$\pm$0.009 & 0.836$\pm$0.005 & 12108 & 87.89$\pm$0.08 \\
    Static Iso.~\cite{kuzucu2024calibration} & 0.759$\pm$0.007 & 0.811$\pm$0.006 & 12803 & 87.96$\pm$0.10 \\
    $g^{\text{cls}}$ only & 0.695$\pm$0.009 & 0.865$\pm$0.009 & 13127 & 87.57$\pm$0.08 \\
    $g^{\text{loc}}$ only & 0.731$\pm$0.013 & 0.846$\pm$0.006 & 12745 & 87.74$\pm$0.01 \\
    Dual ($g^{\text{cls}}{+}g^{\text{loc}}$) & 0.784$\pm$0.006 & 0.817$\pm$0.004 & 11555 & 88.06$\pm$0.10 \\
    CGJS only & 0.751$\pm$0.007 & 0.835$\pm$0.009 & 11851 & 88.04$\pm$0.11 \\
    \textbf{Full CalibPL} & \textbf{0.840}$\pm$0.006 & 0.797$\pm$0.007 & 10494 & \textbf{88.48}$\pm$0.11 \\
    \bottomrule
  \end{tabular}}
  \label{tab:supp_sku_full}
\end{table}

Note: Static Iso.\ (following Kuzucu~\etal) is included as an additional
ablation row. It improves over Temperature Scaling (87.96 vs.\ 87.89 AP$_{50}$)
but falls short of full CalibPL (88.48) because it lacks dynamic re-fitting.

% ============================================================================
\section*{H. Failure Cases and Limitations}
\label{sec:failures}
% ============================================================================

\textbf{Very high overlap ($>0.8$ IoU between GT boxes).}
In pathological cases where two distinct products are nearly identical in
appearance and position, the CGJS gate may incorrectly accept a mis-placed
pseudo-label because both spatial and semantic consistency hold for the wrong
box. This occurs in $\sim$2.1\% of accepted pseudo-labels (measured by
post-hoc IoU matching).

\textbf{Low labeled-data regime ($<$500 labeled boxes).}
Below $\sim$500 labeled boxes, the isotonic regression calibrators overfit to
the calibration set, producing unreliable mappings. The method degrades
gracefully (falls back to fixed threshold behavior) but does not provide
strong calibration benefits in this extreme regime.

\textbf{Anchor-free detectors.}
Our proof assumes score-ranked suppression (NMS). Modern anchor-free detectors
(FCOS, DETR) use different suppression mechanisms. The NMS Tail Amplification
effect still applies qualitatively wherever a maximum-score selection biases
survivors, but the quantitative bound $\Delta_n$ may differ. We leave
anchor-free adaptation to future work.

\end{document}
